\documentclass[preprint, 3p,
authoryear]{elsarticle} %review=doublespace preprint=single 5p=2 column
%%% Begin My package additions %%%%%%%%%%%%%%%%%%%

\usepackage[hyphens]{url}

  \journal{CERC 2026?} % Sets Journal name

\usepackage{graphicx}
%%%%%%%%%%%%%%%% end my additions to header

\usepackage[T1]{fontenc}
\usepackage{lmodern}
\usepackage{amssymb,amsmath}
% TODO: Currently lineno needs to be loaded after amsmath because of conflict
% https://github.com/latex-lineno/lineno/issues/5
\usepackage{lineno} % add
\usepackage{ifxetex,ifluatex}
\usepackage{fixltx2e} % provides \textsubscript
% use upquote if available, for straight quotes in verbatim environments
\IfFileExists{upquote.sty}{\usepackage{upquote}}{}
\ifnum 0\ifxetex 1\fi\ifluatex 1\fi=0 % if pdftex
  \usepackage[utf8]{inputenc}
\else % if luatex or xelatex
  \usepackage{fontspec}
  \ifxetex
    \usepackage{xltxtra,xunicode}
  \fi
  \defaultfontfeatures{Mapping=tex-text,Scale=MatchLowercase}
  \newcommand{\euro}{€}
\fi
% use microtype if available
\IfFileExists{microtype.sty}{\usepackage{microtype}}{}
\usepackage[]{natbib}
\bibliographystyle{plainnat}

\ifxetex
  \usepackage[setpagesize=false, % page size defined by xetex
              unicode=false, % unicode breaks when used with xetex
              xetex]{hyperref}
\else
  \usepackage[unicode=true]{hyperref}
\fi
\hypersetup{breaklinks=true,
            bookmarks=true,
            pdfauthor={},
            pdftitle={Mobilty hub information: practice synthesis and development of tactical urbanism tools},
            colorlinks=false,
            urlcolor=blue,
            linkcolor=magenta,
            pdfborder={0 0 0}}

\setcounter{secnumdepth}{5}
% Pandoc toggle for numbering sections (defaults to be off)


% tightlist command for lists without linebreak
\providecommand{\tightlist}{%
  \setlength{\itemsep}{0pt}\setlength{\parskip}{0pt}}

% From pandoc table feature
\usepackage{longtable,booktabs,array}
\usepackage{calc} % for calculating minipage widths
% Correct order of tables after \paragraph or \subparagraph
\usepackage{etoolbox}
\makeatletter
\patchcmd\longtable{\par}{\if@noskipsec\mbox{}\fi\par}{}{}
\makeatother
% Allow footnotes in longtable head/foot
\IfFileExists{footnotehyper.sty}{\usepackage{footnotehyper}}{\usepackage{footnote}}
\makesavenoteenv{longtable}



\usepackage{subfig}
\usepackage{booktabs}
\usepackage{longtable}
\usepackage{array}
\usepackage{multirow}
\usepackage{wrapfig}
\usepackage{float}
\usepackage{colortbl}
\usepackage{pdflscape}
\usepackage{tabu}
\usepackage{threeparttable}
\usepackage{threeparttablex}
\usepackage[normalem]{ulem}
\usepackage{makecell}
\usepackage{xcolor}



\begin{document}


\begin{frontmatter}

  \title{Mobilty hub information: practice synthesis and development of
tactical urbanism tools}
    \author[]{James Reynolds%
  \corref{cor1}%
  \fnref{1}}
   \ead{julianj.reynolds@atu.ie} 
    \author[]{Brian McCann%
  %
  \fnref{1}}
   \ead{brian.mccann@atu.ie} 
    \author[]{Noreen Armstrong%
  %
  \fnref{2}}
   \ead{noreen.armstrong@atu.ie} 
    \author[]{Sarbast Molsem%
  %
  \fnref{3}}
   \ead{molsem@tcd.ie} 
    \author[]{Agnieszka Stefaniec%
  %
  \fnref{4}}
   \ead{a.stefaniec@soton.ac.uk} 
    \author[]{Brian Caulfield%
  %
  \fnref{3}}
   \ead{brian.caulfield@tcd.ie} 
    \author[]{Amaya Vega%
  %
  \fnref{2}}
   \ead{amaya.Vega@atu.ie} 
      \cortext[cor1]{Corresponding author}
    \fntext[1]{Atlantic Technological University, Sligo campus, Ireland}
    \fntext[2]{Atlantic Technological University, Galway City campus,
Ireland}
    \fntext[3]{Centre for Transport Research, Dept of Civil, Structural
and Environmental Engineering, Trinity College Dublin, Ireland}
    \fntext[4]{Department of Decision Analytics and Risk at Southampton
Business School, University of Southampton, UK}
  
  \begin{abstract}
  
  \end{abstract}
  
 \end{frontmatter}

\section{Introduction}\label{introduction}

Shared mobility hubs are a relatively new development in transport
systems.\\
These provide various shared modes, services and other offerings at a
nearby location, and so help support shifts towards more sustainable and
less-carbon intensive travel.\\
\citet{Roukouni:2023aa} developed a definition from a literature
review\footnote{ \emph{``A shared mobility hub is a place where
  different transport modes are integrated seamlessly, promoting
  efficient and sustainable urban mobility. Emphasis is given on shared
  mobility options such as bikes, scooters and cars. Connection with
  public transport (transit) is desirable but not necessary. The
  integration between mobility suppliers is important to provide a
  seamless, flexible connection at these transfer points. Ensuring an
  enjoyable experience for travelers is a crucial part of the hub
  concept as well, and therefore the hub should be considered not just
  as a transfer node but as a place where a variety of services come
  together, such as logistics services or restaurants.''} (Roukouni et
  al.~2023)}, and this emphasises the need for mobility hubs to be
integrated and to provide seamless transfers from one mode to another. A
similar need to support transfers between services, access and egress
from stations, stops and other facilities has long been a focus for the
design of transit maps and other public information (see for example
\citet{Ovenden:2003aa} and \citet{London-Transport-Museum:2016aa}).

Public information provision will likely be an important consideration
in developing and implementing mobility hubs. However, in a case study
of mobility hubs in London, Bristol and Birmingham
\citet{Reynolds:2026aa} found a range of maps, informational pedestals
and other signage at some hubs, but little to no information at others.
\citet{Reynolds:2026aa} also developed an initial set of software tools
for generating such material for other locations using Open Street Maps
(OSM) data, and the Geographic Information System (GIS) and other
functionality available in the R programming language. However, this
work was based on examples from just Bristol and Birmingham, rather than
from a wider selection of mobility hub informational materials developed
by practitioners and in use in real world contexts.

This paper, therefore, seeks to improve understandings of mobility hub
information provision through review of existing practices and
development of refined versions thorugh synthesis and an iternative
design processes. The study follows the Double Diamond approach
\cite{UK-Design-Council:0aa} by stepping through alternating divergent
and convergent stages of discovery, definition, development and
delivery. Outputs of the research include software tools for producing
mobility hub information products in the style of Bristol, Birmingham
and other cases included in the study. A amalgamated set of designs are
also developed and applied to a shared electric mobility hub site in
Galway, Ireland, as an input to upcoming future research involving
stakeholder co-design and public engagement for user testing purposes.

\section{Research context}\label{research-context}

Origins of the mobility hub concept can be traced back at least to
regional planning activities in Toronto, Canada, in the 2000s
\citep[\citet{Metrolinx:2008aa},
\citet{Salsberg:2010aa}]{Engel-Yan:2012aa}. To a larger extent, however,
railway stations, tram and bus stops, park-and-ride facilities, and many
other locations supporting modal interchange and tavellers are provided
might be considered a `mobility hub'. \citet{Roukouni:2023aa}, however,
notes the emphasis on shared and micro-mobility offerings at mobility
hubs, although perhaps this might similarly broaden the possible
locations given that shared bicycles, e-bikes and e-scooters can now be
relatively easily added to a mix of transport offerings.

Micro-mobility systems, especially of the dockless variety, are not
without their issues. \citet{Meigs:2023aa} and \citet{Zhang:2022aa}
discuss some of the problems relating to battery swapping,
recharging-related fires, inappropriate user behaviours and parking, and
a lack of standards, quality and system management. Shared
\emph{electric} mobility hubs (sometimes `e-Hubs') may have the
potential to reduce some of these, as these incorporate on-site
recharging of vehicles which eliminates many of the issues associated
with operational staff travel for battery swapping and inappropriate
recharging facilities. Power grid connections and the other
infrastructure installed for an electrified hub might also allow other
services to be be provided, for example publicly-accessibility charging
facilities for electric cars.

The \citet{Reynolds:2026aa} study involves one such e-Hub, being a site
in Galway that is part of the ElectRic shared mObility hUbS Trial
(ROBUST) project currently underway in Ireland \cite{ROBUST:2025aa}. As
shown in Figure \ref{fig:CERC_output}, the output of software developed
by \citet{Reynolds:2026aa} generally replicates the informational
displays present at the Gainsborough Square mobility hub in Bristol,
although without the need for bespoke graphic design and mapping.
Rather, the software uses publicaly-available data from Open Street Maps
(OSM) to produce site-specific mobility hub information that might be
sufficient for a pop-up, trial or other tactical-urbanism-style
implementation.

\begin{figure}

{\centering \includegraphics[width=1\linewidth]{graphics/CERC_outputs} 

}

\caption{Gainsborough Square, Bristol, mobility hub signage (left) and mockup for ROBUST Galway site (right). Source: Reynolds et al (in press)}\label{fig:CERC_output}
\end{figure}

While this output is might provide practitioners, researchers and
activists with enough for an initial concept or temporary implementation
of a Bristol-style information display, there appears to be
opportunities for further development. Notably, it is unclear what other
styles of informational displays at mobility hubs have already been
developed, and what might represent best practices in this area. Nor is
it clear what other forms of public information are being used, such as
websites or apps, to promote mobility hubs or to define them as places
and components of larger transportation systems.

The study reported in this paper, therefore, involves cases studies of
further mobility hub locations and a synthesis of current practices. As
dicussed in the following section, the methodology also involves use of
an iterative design process to better understand the design problem that
such information provision measures are seeking to address, various
approaches developed by practitioners, and to deliver a `model'
solution.

\section{Methodology}\label{methodology}

Techniques from industrial design and related fields are readily
applicable to, but perhaps under utilized in, the field of transport
research \citep{Coxon:2018aa,Coxon:2021ab}. Similarly, case-based
research is often under utilized due to a narrow focus on methods that
involve random sampling and a search for statistical significance,
because of confusion about or lack of acceptance of the methodology, or
because of unfamiliarity with the approach (see
\citet{Denscombe:2007aa, Yin:2018ab} and others. Here, the adopted
research methodology combines the Double Diamond design approach, as per
\citet{UK-Design-Council:0aa}, with case studies of various mobility
hubs selected from UK and European locations.

The cases themselves form the primary data source for the discovery
phase of the research. On-site and online data collection was used to
explore information provision at and around mobility hubs, and develop a
definition of the design problem to be addressed in the third
(development) and fourth (delivery) phases. The cases again provided the
primary input for the development phase, which involved the development
of maps, signs, websites and other informational offerings along the
lines of those present at each case location. Finally, a combined `best
practices' version was created and delivered through standard R package
development processes (as per \citet{wickham2023r}).

Hubs were selected for inclusion as cases in the study on the basis of
being particularly unusual, relevatory or leading examples, as per
\citet{Eisenhardt:1989aa, Voss:2002aa, Denscombe:2007aa, Yin:2018ab} and
others. The larger sample available at
https://www.como.org.uk/mobility-hubs/built-and-planned-hubs provided a
pool from which potential cases were drawn.

\section{Results}\label{results}

\subsection{Discovery}\label{discovery}

\begin{longtable}[]{@{}ll@{}}
\caption{Mobility hub components and information offerings. Source:
authors' assessment.}\tabularnewline
\toprule\noalign{}
Component & Bristol \\
\midrule\noalign{}
\endfirsthead
\toprule\noalign{}
Component & Bristol \\
\midrule\noalign{}
\endhead
\bottomrule\noalign{}
\endlastfoot
Bike & Bike \\
eBike & X \\
Cargobike & Cargobike \\
eCargobike & X \\
Scooter & Scooter \\
eScooter & X \\
Bus & X \\
Tram & Tram \\
Train & Train \\
Taxi & Taxi \\
Ride-hail & Ride-hail \\
Share-car & X \\
Share eCar & X \\
Bike parking & X \\
Car parking & Car parking \\
Bike repair & X \\
eBike charge & eBike charge \\
eCar charge & X \\
Toilets & X \\
Phone charge & X \\
Wifi & X \\
Shelter & Shelter \\
Seating & X, \\
Water & Water \\
Food & X \\
Groceries & X \\
Library & X \\
Pedestal & X \\
Local map & X \\
Area map & X \\
Visual indicators & X \\
Website & Website \\
App & App \\
Pricing & Pricing \\
\end{longtable}

\begin{figure}

{\centering \includegraphics[width=1\linewidth]{graphics/andrew_road} 

}

\caption{Andrew Road hub, selected images. Source: authors.}\label{fig:andrew_road_images}
\end{figure}

\subsection{Definition}\label{definition}

\subsection{Development}\label{development}

\subsection{Delivery}\label{delivery}

\section{Discussion}\label{discussion}

\section{Conclusions}\label{conclusions}

\begin{longtable}[]{@{}llllllll@{}}
\caption{Author contribution matrix, as per CRediT (Contribution Roles
Taxonomy)}\tabularnewline
\toprule\noalign{}
Component & NA & BC & BM & SM & JR & AS & AV \\
\midrule\noalign{}
\endfirsthead
\toprule\noalign{}
Component & NA & BC & BM & SM & JR & AS & AV \\
\midrule\noalign{}
\endhead
\bottomrule\noalign{}
\endlastfoot
Conceptualization & & & & & X & & \\
Data Curation & & & & & X & & \\
Formal Analysis & & & & & X & & \\
Funding Acquisition & & X & & & & X & \\
Investigation & & & & & X & & \\
Methodology & & & & & X & & \\
Project Administration & & & & X & & & \\
Resources & & & & & & & \\
Software & & & & & X & & \\
Supervision & & & X & & & & X \\
Validation & & & & & X & & \\
Visualization & & & & & X & & \\
Writing -- Original Draft Preparation & & & & & X & & \\
Writing -- Review \& Editing & & & X & & X & & \\
\end{longtable}

\renewcommand\refname{References}
\bibliography{References.bib, packages.bib}


\end{document}
